\section{Methods}
\label{sec:methods}

\subsection{What we are predicting}

The data come from ClinicalTrials.gov through the HINT benchmark, with further
processing by the course staff~\cite{fu2022hint,clinicaltrials}. There are 1,561
trials that stopped early. The training file has 1,171 of them. Only 157 have
visible reason labels: the earliest 10\% by start year, plus 40 extra trials
chosen so the rare reasons are not missing. The other 1,014 training trials have
every reason marked $-100$. That value means ``unknown,'' not ``this was not a
reason.'' The validation file has 78 fully labeled trials. The test file has 312
later trials, from 2010 to 2016, with no labels.

A trial can stop for more than one reason. The five labels are Safety, Funding,
Efficacy, Recruitment, and Other. Among the 235 trials with real labels,
Recruitment appears on 64\%, Other on 15\%, Efficacy on 13\%, Safety on 11\%, and
Funding on 5\%. The labeled trials are mostly early, and the test trials are the
latest, so these rates will not match the test set. Recruitment is more common
early, and Safety is more common later.

The score we care about is sample-wise F1: F1 computed for each trial, then
averaged over trials. We also report F1 for each reason and the macro average,
which is the plain average of those five F1 scores. One validation trial moves
sample-wise F1 by about 1.3 points, so a small gap on 78 trials is weak evidence.
Every comparison below was therefore checked a second way, with the same saved
grouped cross-validation: eight repeats of five folds on the 157 training labels,
and the same plan on all 235 labeled trials. We also scored three time-ordered
windows, 98 held-out trials in total. Where a single validation number and the
repeated checks disagree, we follow the repeated checks.

\subsection{Which inputs we kept}

Each trial has a 150-number embedding of its stopping-reason text. The original
text is not in the files. The embedding was reduced and randomly rotated, so a
single coordinate has no meaning, and the values are small. We standardize them
with a scaler fit only on the rows the current split is allowed to learn from.
Trial IDs are used to line up rows. They are never used as inputs.

We tried the other fields and then left them out. Start year, eligibility text,
disease names, drug pathways, and Morgan fingerprints of the drug
structure~\cite{rogers2010extended,rdkit,milacic2024reactome} did not raise the
validation score. The reason is simple. The embedding is already a summary of the
text the labels were written from. The other fields describe the trial, not why
it stopped. A trial can run out of money in any disease. Adding them also means
adding many new inputs to a model that starts with only 157 real labels. About
44\% of trials have no matched drug, so a fingerprint block is all zeros for
those rows. Year is a special trap: it fits the early labeled trials and then
has to predict later ones. Section~\ref{sec:results} gives the numbers.

\subsection{How missing labels are created}

The assignment asks for at least one unsupervised method to fill in missing
labels. We tried three. $k$-means and DBSCAN were useful checks. A careful
DBSCAN labeling, followed by logistic regression, reached about 0.910 sample-wise
F1 on the 78 validation trials in an early comparison. We did not submit it.
That number came from a different training setup, and it did not win the later
grouped checks against the model below. On its own, clustering mostly separates
Recruitment from everything else, because Recruitment dominates the data.

What we kept is a Gaussian mixture~\cite{dempster1977em,pedregosa2011scikit}.
After standardization we keep 30 principal components, about 77\% of the
variance, and fit 30 diagonal-covariance components. The components are started
from $k$-means++ assignments. We run three restarts from
\texttt{random\_state}${}=0$ and keep the restart with the best log-likelihood.
Expectation--maximization stops when the average log-likelihood improves by less
than $10^{-3}$, or after 100 iterations. The kept run converged in 29
iterations. The mixture is fit only on rows the current split may see. Test rows
are not included.

Each component is turned into a reason profile by averaging the real labels of
the trials assigned to it, weighted by how strongly the mixture assigns each
trial. A trial with two reasons contributes half a count to each, so every trial
has the same total weight. We add one pseudo-count of the overall reason rates.
That stops a component with few labeled trials from collapsing to a single
reason. The mixture's guess for a new trial is the average of these profiles,
weighted by the trial's component memberships.

We do not trust that guess as a label. In the teaching step it has only a 5\%
vote. The classifier blend below has the other 95\%.

\subsection{The classifier}

The submitted classifier is one logistic regression over the five reasons
($C = 0.1$, at most 3,000 iterations, tolerance $10^{-4}$), not five separate
classifiers. A trial with two reasons is entered once for each reason, and one
unit of weight is split evenly between those copies. The class weights are
computed from the real labels only, then frozen. Rare reasons are upweighted so
Recruitment does not drown them out. A label the model made up later cannot
change those weights. A made-up row also counts as only a quarter of a real row.

We average this logistic model with two others, using fixed weights. Cosine
nearest neighbors get 30\%. Shrinkage linear discriminant analysis gets 20\%.
Logistic regression gets the remaining 50\%. The neighbor model looks at the 10
closest trials and weights them by the square of the cosine similarity, so a
close match counts for much more than a distant one. A neighbor that has two
reasons votes for both. The discriminant analysis uses automatic shrinkage and
equal class priors, so it is not pulled toward Recruitment just because that
reason is common. All three models see the full 150 standardized embeddings.
Only the mixture uses the 30 principal components.

This blend is the later version of our first strong model, with three changes
that the grouped checks supported. We froze the class weights on the real
labels, we let neighbors vote for every reason they carry, and we cut the
mixture's teaching vote from 20\% to 5\%.

\subsection{Teaching the model, then deciding}

We add unlabeled trials in two rounds~\cite{yarowsky1995unsupervised,triguero2015selflabeled}.
Each round scores every still-unlabeled trial with the 95/5 blend above. If its
highest reason scores at least 0.8, we treat the decision rule below as that
trial's label, add the row at quarter weight, and refit. Each round starts again
from the real labels plus the rows accepted in that round. It does not keep
stacking old guesses. In the final fit, on all 235 real labels, the first round
accepts 685 unlabeled trials and the second accepts 757. Their total weight is
about 103 and then 120, against 235 for the real labels, so the guesses cannot
outvote the data.

The course also asks for a reason label on every originally masked training row.
After the final fit we predict all 1,014 of those rows. That file is not the
same thing as the 757 rows used during training. The 157 original training
labels are copied through unchanged.

Every trial in this corpus stopped for a recorded reason, so a prediction of
``no reason'' is always wrong. We always keep the highest-scoring reason, and we
add any other reason scoring at least 0.3. These scores are for making that
choice. They are not calibrated probabilities.

\subsection{What the final fit is allowed to see}

Model choice used the 78 validation trials, with the model trained on the 157
labeled training trials. The submitted fit is different. It retrains the same
fixed recipe on all 235 real labels, and it may use the 1,014 masked training
rows for the mixture and for generated labels. The 312 test rows are used only
to make the final predictions. Because the validation labels are inside that
final fit, the 0.891 validation score is not an independent score of the
submitted file.

Our first submitted model, at a reported Kaggle score of 0.89423, had built its
scaler and clusters with the test rows included. We stopped doing that. The
later models, including the one below, never see a test row until prediction.
